% Authoritative LaTeX chapter. Edit this file for future revisions.
% Initially migrated 2026-09-11; no Markdown import occurs during PDF builds.
\chapter{기본 개념과 관측 기하}
\label{chap:01_concepts_geometry}
매뉴얼 목차 · \hyperref[chap:05_output_formats]{출력 포맷}

이 문서는 \texttt{code/\allowbreak{}OCRT\_\allowbreak{}C/\allowbreak{}src}의 공개 입력과 계산 결과를 기준으로 한다. 확인 기준은 Git 커밋 \texttt{54861c68e35ec5aa\allowbreak{}a3938fad53e93dc2\allowbreak{}dfa7eb1d}이며, 실행 파일의 버전 문자열만으로 최신 소스 여부를 판정하지 않는다. 과거 검증 보고서는 당시 버전에 대한 기록이다.

\section{OCRT가 계산하는 것}
\label{sec:auto:01_concepts_geometry:1}
OCRT는 태양광이 대기, 해수면, 물속에서 산란·흡수·반사·투과되는 과정을 계산하는 복사전달 모델이다. 같은 파장의 탄성 산란에 대해 Stokes \texttt{I}, \texttt{Q}, \texttt{U}를 계산한다. \texttt{I}는 총 복사휘도에 해당하고, \texttt{Q}, \texttt{U}는 선택한 기준면에 대한 선편광 성분이다. 원편광 \texttt{V}는 공개 출력에 없다.

입력은 관측 기하, 파장, 대기압·에어로졸·기체, 해수면 바람, 물의 광학 특성이다. 주요 출력은 다음과 같다.


{\TableFont
\begin{longtable}{@{}>{\raggedright\arraybackslash}p{\dimexpr 0.3400\linewidth-4.00pt\relax}>{\raggedright\arraybackslash}p{\dimexpr 0.6600\linewidth-4.00pt\relax}@{}}
\toprule
\textbf{위치·수량} & \textbf{의미} \\
\midrule
\endfirsthead
\toprule
\textbf{위치·수량} & \textbf{의미} \\
\midrule
\endhead
\bottomrule
\endfoot
TOA & 대기 상단에서 위성 방향으로 나가는 신호 \\
\texttt{0+} & 수면 바로 위, 공기 쪽 \\
\texttt{0-\allowbreak{}} & 수면 바로 아래, 물 쪽 \\
\texttt{rho} & 태양 입사량으로 정규화한 방향별 TOA 반사도, 무차원 \\
\texttt{Rrs} & \texttt{Lu(0+)/\allowbreak{}Ed(0+)}, 수면 위 원격탐사반사도, sr\textsuperscript{-}\textsuperscript{1} \\
\texttt{rrs} & \texttt{Lu(0-\allowbreak{})/\allowbreak{}Ed(0-\allowbreak{})}, 수면 아래 원격탐사반사도, sr\textsuperscript{-}\textsuperscript{1} \\
\texttt{Lu}, \texttt{Ed}, \texttt{Eu} & 각각 상향 복사휘도, 하향 복사조도, 상향 복사조도 \\
\end{longtable}
}

\texttt{Rrs}와 \texttt{rho}는 분모와 단위가 다르다. \texttt{Rrs}에 \ensuremath{\pi}를 곱한 수면 반사도도 대기 경로가 포함된 TOA \texttt{rho}와 같지 않다. 원시 \texttt{Lu}, \texttt{Ed}는 모델별 입사광 정규화가 적용된 값이므로 \hyperref[chap:05_output_formats]{출력의 정규화 규칙}을 먼저 확인한다.

계산은 수평으로 균일한 층을 가정한다. 다중산란은 SOS(Successive Orders of Scattering, 산란 차수별 누적)로 계산하고, 방위각 의존성은 Fourier 모드로 표현한다. \texttt{-\allowbreak{}-\allowbreak{}pssa}를 켜면 현재 구현은 IPSS 구면 보정을 적용한다. 완전한 3차원 지형·구름 복사전달이나 임의 해저 지형을 입력하는 모델로 해석하지 않는다.

\section{표면 모드를 먼저 선택한다}
\label{sec:auto:01_concepts_geometry:2}

{\TableFont
\begin{longtable}{@{}>{\raggedright\arraybackslash}p{\dimexpr 0.3000\linewidth-5.33pt\relax}>{\raggedright\arraybackslash}p{\dimexpr 0.2300\linewidth-5.33pt\relax}>{\raggedright\arraybackslash}p{\dimexpr 0.4700\linewidth-5.33pt\relax}@{}}
\toprule
\textbf{\texttt{-\allowbreak{}-\allowbreak{}surface}} & \textbf{계산 범위} & \textbf{대표 용도} \\
\midrule
\endfirsthead
\toprule
\textbf{\texttt{-\allowbreak{}-\allowbreak{}surface}} & \textbf{계산 범위} & \textbf{대표 용도} \\
\midrule
\endhead
\bottomrule
\endfoot
\texttt{black} & 반사하지 않는 하부 경계 위의 대기 & 순수 대기 경로, Rayleigh 산란 확인 \\
\texttt{flat} & 평탄한 Fresnel 반사 경계 위의 대기 & 평탄면 Fresnel 기준 계산 \\
\texttt{black\_\allowbreak{}fresnel\_\allowbreak{}ocean} & 바람에 따른 Fresnel 해수면과 대기, 물속 수출광 없음 & 해양 대기보정용 Rayleigh/에어로졸 경로 LUT \\
\texttt{ocean} & 대기--해수면--수중 복사전달의 결합 & \texttt{Rrs}, \texttt{rrs}, TOA 해양 신호·모의자료 \\
\end{longtable}
}

\texttt{ocean}에서는 물의 특성을 \texttt{-\allowbreak{}-\allowbreak{}water-\allowbreak{}model ocrt}, \texttt{ccrr}, \texttt{iop} 중 하나로 지정한다. \texttt{ocrt}와 \texttt{ccrr}는 구성성분을 광학 특성으로 바꾸는 방법을 선택하고, \texttt{iop}는 흡수·산란·후방산란 계수와 위상함수를 직접 공급한다. 같은 이름의 농도라도 구성성분 모델에 따라 광학 특성이 달라질 수 있다.

IOP(Inherent Optical Properties, 고유광학특성)의 \texttt{a}, \texttt{b}, \texttt{bb}는 각각 흡수, 산란, 후방산란 계수이며 단위는 m\textsuperscript{-}\textsuperscript{1}이다. 소광 계수는 \texttt{c=\allowbreak{}a+b}, 단일산란 알베도는 \texttt{omega=\allowbreak{}b/\allowbreak{}(a+b)}다. \texttt{bb}는 \texttt{b}의 후방 반구 성분이므로 \texttt{a+b+bb}를 소광 계수로 사용하지 않는다.

\section{SZA, VZA, RAA: 반드시 지상 화소 기준}
\label{sec:auto:01_concepts_geometry:3}

{\TableFont
\begin{longtable}{@{}>{\raggedright\arraybackslash}p{\dimexpr 0.3000\linewidth-5.33pt\relax}>{\raggedright\arraybackslash}p{\dimexpr 0.2300\linewidth-5.33pt\relax}>{\raggedright\arraybackslash}p{\dimexpr 0.4700\linewidth-5.33pt\relax}@{}}
\toprule
\textbf{입력·CSV 이름} & \textbf{정의} & \textbf{단위} \\
\midrule
\endfirsthead
\toprule
\textbf{입력·CSV 이름} & \textbf{정의} & \textbf{단위} \\
\midrule
\endhead
\bottomrule
\endfoot
\texttt{-\allowbreak{}-\allowbreak{}sza}, \texttt{sza\_\allowbreak{}deg} & 지상 화소의 위쪽 수직선과 태양을 향하는 방향 사이의 각도 & 도 \\
\texttt{-\allowbreak{}-\allowbreak{}vza}, \texttt{vza\_\allowbreak{}deg} & 지상 화소의 위쪽 수직선과 센서를 향하는 방향 사이의 각도 & 도 \\
\texttt{-\allowbreak{}-\allowbreak{}raa}, \texttt{raa\_\allowbreak{}deg} & 아래 운영 정의를 따르는 OCRT 공개 상대방위각 & 도 \\
\end{longtable}
}

\texttt{VZA=\allowbreak{}0\textdegree{}}는 천저 관측이다. VZA는 위성에서 내려다본 off-nadir angle과 구별한다. 구면 기하에서는 두 각이 다를 수 있다. \texttt{-\allowbreak{}-\allowbreak{}pssa}를 사용해도 입력은 지상 화소 기준이다. \texttt{ocean}에서 VZA는 공기 중 관측각이며, 표준 경로는 물속 각도를 내부 굴절 변환으로 구한다. 동일한 VZA 숫자를 수중 각도로 다시 해석하면 안 된다.

낮의 상향 관측에서는 일반적으로 \texttt{0 \ensuremath{\leq} SZA < 90}, \texttt{0 \ensuremath{\leq} VZA < 90}의 유효 기하를 준비한다. 기본 LUT는 VZA \texttt{0--85\textdegree{}}이며 정확한 격자는 출력 CSV의 좌표를 확인한다. RAA는 \texttt{[0,\allowbreak{}360)}으로 정규화해 입력하는 것을 권장한다. \texttt{360\textdegree{}}는 물리적으로 \texttt{0\textdegree{}}와 같은 방향이며 LUT에는 중복 끝점으로 넣지 않는다.

그림~\ref{fig:geometry-zenith-refraction}은 수직선에서 재는 두 천정각과 수중 굴절을 구분한다. 그림~\ref{fig:geometry-curvature}은 왜 위성의 off-nadir 각도를 \texttt{--vza}에 그대로 넣으면 안 되는지 보여준다.

% Editable vector artwork. Both editions use the same physical geometry.
\begin{figure}[H]
\centering
\begin{tikzpicture}[x=1cm,y=1cm,font=\sffamily\fontsize{9}{11}\selectfont,
  ray/.style={line width=1.05pt,-{Latex[length=2.2mm]}},
  axis/.style={ocrtGray,dashed,line width=.55pt},
  angle/.style={ocrtNavy,line width=.8pt},
  note/.style={align=center,text width=6.8cm}]
\path[use as bounding box] (0,-3.6) rectangle (15.8,3.9);
\begin{scope}[shift={(3.6,0)}]
  \node[anchor=north,font=\sffamily\bfseries] at (0,3.9)
    {\FigLang{(a) 화소에서 위쪽을 바라본 각도}{(a) Look directions from the pixel}};
  \fill[ocrtTeal!8] (-3.3,-.55) rectangle (3.3,0);
  \draw[ocrtTeal,line width=.8pt] (-3.3,0)--(3.3,0);
  \draw[axis] (0,-.35)--(0,2.75);
  \node[anchor=south,fill=white,inner sep=1.5pt] at (0,2.75)
    {\FigLang{국소 수직선}{Local vertical}};
  \draw[ocrtGray,dashed,-{Latex[length=2mm]}] (0,0)--(125:2.65);
  \draw[ray,orange!85!black] (125:2.25)--(125:.35);
  \draw[ray,ocrtTeal] (0,0)--(45:2.8);
  \node[anchor=south,align=center] at (-1.85,2.37)
    {\FigLang{태양을 향한 방향}{Toward the Sun}};
  \node[anchor=south,align=center] at (2.02,2.08)
    {\FigLang{센서를 향한 방향}{Toward the sensor}};
  \draw[angle] (90:.85) arc[start angle=90,end angle=125,radius=.85];
  \node[ocrtNavy] at (-.52,1.14) {SZA};
  \draw[angle] (45:1.12) arc[start angle=45,end angle=90,radius=1.12];
  \node[ocrtNavy] at (.66,1.37) {VZA};
  \fill[ocrtNavy] (0,0) circle (2pt);
  \node[anchor=north] at (0,-.12) {\FigLang{지상 화소}{Ground pixel}};
  \node[note,anchor=north] at (0,-.95)
    {\FigLang{태양을 바라보는 방향과 입사광의 진행 방향은 서로 반대다.}{The Sunward look direction is opposite to incoming photon travel.}};
  \draw[ray,orange!85!black] (-2.75,-2.13)--(-1.9,-2.13);
  \node[anchor=west] at (-1.7,-2.13) {\FigLang{입사 태양광}{Incoming sunlight}};
  \draw[ray,ocrtTeal] (-2.75,-2.68)--(-1.9,-2.68);
  \node[anchor=west] at (-1.7,-2.68) {\FigLang{센서로 나가는 빛}{Light toward the sensor}};
\end{scope}
\draw[ocrtLine] (7.8,-3.35)--(7.8,3.55);
\begin{scope}[shift={(11.8,0)}]
  \node[anchor=north,font=\sffamily\bfseries] at (0,3.9)
    {\FigLang{(b) 공기 중 VZA와 수중 각도}{(b) Air VZA and underwater angle}};
  \fill[ocrtTeal!8] (-3.4,-1.9) rectangle (3.4,0);
  \draw[ocrtTeal,line width=.8pt] (-3.4,0)--(3.4,0);
  \draw[axis] (0,-1.8)--(0,2.65);
  \draw[ray,ocrtTeal] (-.847,-1.55)--(0,0);
  \draw[ray,ocrtTeal] (0,0)--(50:2.65);
  \node[anchor=south,align=center] at (1.9,2.1)
    {\FigLang{센서}{Sensor}};
  \node[anchor=west] at (-3.18,.42) {\FigLang{공기}{Air}};
  \node[anchor=west] at (-3.18,-.43) {\FigLang{물}{Water}};
  \draw[angle] (50:.9) arc[start angle=50,end angle=90,radius=.9];
  \node[ocrtNavy] at (.7,1.17) {$\mathrm{VZA}_{\rm air}$};
  \draw[angle] (-90:.7) arc[start angle=-90,end angle=-118.66,radius=.7];
  \node[ocrtNavy] at (-.63,-1.05) {$\theta_w$};
  \fill[ocrtNavy] (0,0) circle (2pt);
  \node[note,anchor=north] at (0,-2.16)
    {$n_{\rm air}\sin(\mathrm{VZA}_{\rm air})=n_w\sin\theta_w$};
  \node[note,anchor=north] at (0,-2.7)
    {\FigLang{입력: 공기 중 VZA\\수중 각도는 내부에서 계산}{Input: air VZA\\Underwater angle is computed internally}};
\end{scope}
\end{tikzpicture}
\caption{\FigLang{SZA와 VZA는 화소의 위쪽 수직선에서 잰다. (a)는 두 look 방향을 한 단면에 배치한 예시이며 일반 기하의 RAA는 별도로 지정한다. (b)는 평탄한 경계의 굴절 도식이다. 표준 해양 경로의 VZA 입력을 수중 각도로 바꾸어 넣지 않는다.}{SZA and VZA are measured from the upward vertical at the pixel. Panel (a) places both look directions in one cross-section for illustration; RAA is specified separately for general geometry. Panel (b) shows refraction at a flat interface. Do not replace the standard ocean-path air-VZA input with an underwater angle.}}
\label{fig:geometry-zenith-refraction}
\end{figure}

% Angles drawn from exact ray endpoints; exaggerated orbital height is schematic.
\begin{figure}[H]
\centering
\begin{tikzpicture}[x=1cm,y=1cm,font=\sffamily\fontsize{9}{11}\selectfont,
  note/.style={align=left,text width=5.4cm}]
\path[use as bounding box] (0,-3.15) rectangle (15.8,4.72);
\coordinate (C) at (3.6,-2.5);
\coordinate (P) at (3.6,1);
\coordinate (S) at (8.6,4);
\fill[ocrtTeal!9] (C)--($(C)+(25:3.5)$)
  arc[start angle=25,end angle=155,radius=3.5]--cycle;
\draw[ocrtTeal,line width=1pt] ($(C)+(25:3.5)$)
  arc[start angle=25,end angle=155,radius=3.5];
\draw[ocrtLine,dotted,line width=.8pt] ($(C)+(25:3.85)$)
  arc[start angle=25,end angle=155,radius=3.85];
\draw[ocrtGray,dashed] (C)--(3.6,4.25);
\draw[ocrtGray,dashed] (S)--(C);
\draw[ocrtTeal,line width=1.15pt,-{Latex[length=2.2mm]}] (P)--(S);
\draw[ocrtNavy,line width=.8pt] ($(P)+(30.96376:1)$)
  arc[start angle=30.96376,end angle=90,radius=1];
\node[ocrtNavy,fill=white,inner sep=1pt] at (4.24,2.36) {VZA};
\draw[purple!70!black,line width=.8pt] ($(S)+(-149.03624:1.1)$)
  arc[start angle=-149.03624,end angle=-127.56859,radius=1.1];
\node[purple!70!black,fill=white,inner sep=1pt] at (7.25,2.87) {$\beta$};
\fill[ocrtNavy] (P) circle (2pt);
\fill[ocrtNavy] (C) circle (2pt);
\fill[ocrtTeal] (S) circle (2.3pt);
\node[anchor=south,align=center] at (3.6,4.29)
  {\FigLang{화소의 국소 수직선}{Local vertical at the pixel}};
\node[anchor=south,align=center] at (8.45,4.19)
  {\FigLang{위성}{Satellite}};
\node[align=center] at (2.3,1.95)
  {\FigLang{지상 화소}{Ground pixel}};
\draw[ocrtGray,line width=.5pt] (2.8,1.66)--(3.5,1.1);
\node[anchor=north,align=center] at (3.6,-2.69)
  {\FigLang{지구 중심}{Earth center}};
\node[align=center] at (1.3,.95)
  {\FigLang{대기층}{Atmosphere}};
\draw[ocrtGray,line width=.5pt] (2.22,1.02)--(2.5,1.189);
\node[fill=ocrtTeal!9,inner sep=2pt] at (3.22,-.65) {$R$};
\node[rotate=52.43,fill=white,inner sep=2pt] at (5.8,.47) {$R+h$};
\node[anchor=north west,note,font=\sffamily\bfseries] at (10.0,3.9)
  {\FigLang{입력하는 각도는 VZA}{Enter VZA at the pixel}};
\node[anchor=north west,note] at (10.0,3.1)
  {\FigLang{\texttt{--vza}: 화소의 위쪽 수직선과 센서 look 방향 사이의 각도.}{\texttt{--vza}: angle between the pixel's upward vertical and its sensor look direction.}};
\node[anchor=north west,note] at (10.0,1.65)
  {\FigLang{$\beta$: 위성에서 지구 중심을 향하는 천저 방향과 화소를 향하는 방향 사이의 off-nadir 각도.}{$\beta$: satellite off-nadir angle, between the Earth-center direction and the look toward the pixel.}};
\node[anchor=north west,note] at (10.0,-.05)
  {$R\sin(\mathrm{VZA})=(R+h)\sin\beta$};
\node[anchor=north west,note] at (10.0,-.76)
  {\FigLang{$R$: 구의 반지름, $h$: 센서 고도. 직선 광선·구형 지구에서의 기하 관계다.}{$R$: sphere radius; $h$: sensor altitude. This relation assumes straight rays and a spherical Earth.}};
\node[anchor=north west,note] at (10.0,-2.24)
  {\FigLang{\texttt{--pssa}에서도 입력 기준은 화소다.}{The input remains pixel-referenced with \texttt{--pssa}.}};
\end{tikzpicture}
\caption{\FigLang{구면 기하에서 지상 화소의 VZA와 위성의 off-nadir 각도는 일반적으로 다르다. 대기층 두께와 센서 고도는 구별이 쉽도록 과장했으며 비례도·계산 결과가 아니다. 현재 \texttt{--pssa}의 IPSS 보정도 지상 화소 기준 기하를 사용한다. 이 도식은 완전한 3차원 구면 다중산란 모델을 의미하지 않는다.}{In spherical geometry, ground-pixel VZA and satellite off-nadir angle generally differ. Atmospheric thickness and sensor altitude are exaggerated for clarity; this is not a scale drawing or a computed result. The current IPSS correction enabled by \texttt{--pssa} also uses ground-pixel geometry. The schematic does not imply a fully three-dimensional spherical multiple-scattering model.}}
\label{fig:geometry-curvature}
\end{figure}


\section{OCRT RAA의 운영 정의}
\label{sec:auto:01_concepts_geometry:4}
\textbf{OCRT의 직접 선글린트 방향은 \texttt{RAA=\allowbreak{}180\textdegree{}}다.} 평탄면 정반사 기하에서 \texttt{SZA=\allowbreak{}VZA}일 때 이 방향이 정반사 가지에 해당한다. 바람이 있으면 Cox--Munk 경사 분포가 정반사 신호를 주변 각도로 넓힌다.


{\TableFont
\begin{longtable}{@{}>{\raggedright\arraybackslash}p{\dimexpr 0.3400\linewidth-4.00pt\relax}>{\raggedright\arraybackslash}p{\dimexpr 0.6600\linewidth-4.00pt\relax}@{}}
\toprule
\textbf{OCRT RAA} & \textbf{해석} \\
\midrule
\endfirsthead
\toprule
\textbf{OCRT RAA} & \textbf{해석} \\
\midrule
\endhead
\bottomrule
\endfoot
\texttt{0\textdegree{}} & 직접 글린트의 반대 주평면 가지 \\
\texttt{90\textdegree{}} & 주평면에 수직인 한쪽 방위 \\
\texttt{180\textdegree{}} & 직접 글린트 주평면 가지; \texttt{SZA=\allowbreak{}VZA}에서 정반사 기하 \\
\texttt{270\textdegree{}} & \texttt{90\textdegree{}}의 거울 방위; \texttt{I/\allowbreak{}Q}는 대칭이고 \texttt{U}는 반대 부호가 될 수 있음 \\
\texttt{360\textdegree{}} & \texttt{0\textdegree{}}와 동일한 방향 \\
\end{longtable}
}

“태양과 같은 쪽”, “전방”, “후방”이라는 말만으로 RAA를 결정하지 않는다. 태양을 \textbf{바라보는 방향}과 태양광이 \textbf{진행하는 방향}은 수평 방위가 180\textdegree{} 다르기 때문이다. 다른 복사전달 코드나 LUT에서 \texttt{RAA=\allowbreak{}0\textdegree{}}가 글린트라면 그 숫자를 OCRT에 그대로 넣을 수 없다.

그림~\ref{fig:geometry-glint-branches}에서 태양광의 진행 방향과 센서 look 방향을 함께 따라가면 두 가지를 구별하기 쉽다. 그림의 \texttt{RAA=180\textdegree{}}만으로 정반사 중심이 결정되는 것은 아니며 \texttt{SZA=VZA}도 필요하다.

% Equal 30-degree zenith angles. Solar-look and propagation vectors differ.
\begin{figure}[H]
\centering
\begin{tikzpicture}[x=1cm,y=1cm,font=\sffamily\fontsize{9}{11}\selectfont,
  ray/.style={line width=1.15pt,-{Latex[length=2.2mm]}},
  note/.style={align=center,text width=6.9cm}]
\path[use as bounding box] (0,-3.15) rectangle (15.8,4.25);
\begin{scope}[shift={(3.6,0)}]
  \node[anchor=north,font=\sffamily\bfseries] at (0,4.25)
    {(a) RAA $=0^\circ$};
  \fill[ocrtTeal!8] (-3.3,-.4) rectangle (3.3,0);
  \draw[ocrtTeal,line width=.8pt] (-3.3,0)--(3.3,0);
  \draw[ocrtGray,dashed] (0,0)--(0,2.45);
  \draw[ocrtGray,dashed] (0,0)--(120:3.0);
  \draw[ray,orange!85!black] (120:2.65)--(120:1.7);
  \draw[ray,ocrtTeal] (120:.15)--(120:1.35);
  \node[anchor=south,align=center,text width=4.8cm] at (-.9,2.75)
    {\FigLang{태양과 센서의 look 방향 일치}{Solar and sensor look directions coincide}};
  \draw[ocrtNavy] (90:.8) arc[start angle=90,end angle=120,radius=.8];
  \node[anchor=west,align=left] at (.3,1.25)
    {SZA $=$ VZA $=30^\circ$};
  \fill[ocrtNavy] (0,0) circle (2pt);
  \node[anchor=north] at (0,-.12) {\FigLang{화소}{Pixel}};
  \node[note,anchor=north] at (0,-.79)
    {\FigLang{평탄면의 직접 글린트 가지가 아님}{Outside the flat-surface direct-glint branch}};
  \node[note,anchor=north] at (0,-1.88)
    {$\Theta=180^\circ$\\\FigLang{대기 입자의 정확 후방산란 기하}{Exact atmospheric backscattering geometry}};
\end{scope}
\draw[ocrtLine] (7.8,-2.95)--(7.8,3.85);
\begin{scope}[shift={(11.8,0)}]
  \node[anchor=north,font=\sffamily\bfseries] at (0,4.25)
    {(b) RAA $=180^\circ$};
  \fill[ocrtTeal!8] (-3.3,-.4) rectangle (3.3,0);
  \draw[ocrtTeal,line width=.8pt] (-3.3,0)--(3.3,0);
  \draw[ocrtGray,dashed] (0,0)--(0,2.85);
  \draw[ray,orange!85!black] (120:2.95)--(0,0);
  \draw[ray,ocrtTeal] (0,0)--(60:2.95);
  \node[anchor=south] at (-1.5,2.7) {\FigLang{태양}{Sun}};
  \node[anchor=south] at (1.5,2.7) {\FigLang{센서}{Sensor}};
  \draw[ocrtNavy] (90:.9) arc[start angle=90,end angle=120,radius=.9];
  \draw[ocrtNavy] (60:.9) arc[start angle=60,end angle=90,radius=.9];
  \node[ocrtNavy] at (-.55,1.16) {SZA};
  \node[ocrtNavy] at (.55,1.16) {VZA};
  \fill[ocrtNavy] (0,0) circle (2pt);
  \node[anchor=north] at (0,-.12) {\FigLang{화소}{Pixel}};
  \node[note,anchor=north] at (0,-.79)
    {SZA $=$ VZA $=30^\circ$\\\FigLang{평탄면 정반사}{Flat-surface specular reflection}};
  \node[note,anchor=north] at (0,-1.88)
    {$\Theta=120^\circ$\\\FigLang{수면 정반사와 입자 후방산란은 구별}{Surface specular reflection differs from particle backscatter}};
\end{scope}
\end{tikzpicture}
\caption{\FigLang{두 주평면 가지를 같은 천정각으로 비교한 단면도. 주황 화살표는 입사광, 청록 화살표는 센서 방향으로 나가는 빛이다. (a)의 두 화살표는 같은 직선 위에서 서로 반대 방향을 가리킨다. $\Theta$는 입사 태양광의 진행 방향과 상향 관측 방향 사이의 대기 산란각이며 수면 법선에서 잰 각도가 아니다. 바람이 있는 해수면의 글린트는 (b)의 중심 주변으로 넓어질 수 있다.}{Cross-sections comparing the two principal-plane branches at equal zenith angles. Orange arrows show incoming photons and teal arrows show light toward the sensor. The arrows in (a) point in opposite directions on the same line. $\Theta$ is the atmospheric scattering angle between incoming solar propagation and upward viewing, not an angle measured from the surface normal. A wind-roughened surface can broaden glint around the center shown in (b).}}
\label{fig:geometry-glint-branches}
\end{figure}


\subsection{위성 SAA/VAA로부터 계산}
\label{sec:auto:01_concepts_geometry:4.1}
먼저 제품 메타데이터에서 다음을 확인한다.

\begin{enumerate}
\item SAA가 지상 화소에서 \textbf{태양을 향하는} 방위각인지 확인한다.
\item VAA가 지상 화소에서 \textbf{센서를 향하는} 방위각인지 확인한다. 센서에서 지표를 향하는 방위라면 \texttt{180\textdegree{}}를 더해 지상 화소 기준으로 바꾼다.
\item SAA와 VAA의 영점 및 양의 회전 방향을 일치시킨다. 이미 상대방위각으로 제공된 값에 이 절차를 다시 적용하지 않는다.
\end{enumerate}

두 값이 위의 look-azimuth 정의를 만족하고 같은 회전 방향을 쓸 때, 다음은 OCRT의 \texttt{0/\allowbreak{}180\textdegree{}} 가지와 일치하는 명시적인 변환 규약이다.


\begin{ManualCode}
def ocrt_raa_from_look_azimuths(saa_deg, vaa_deg):
    return (vaa_deg - saa_deg) % 360.0
\end{ManualCode}


{\TableFont
\begin{longtable}{@{}>{\raggedright\arraybackslash}p{\dimexpr 0.2300\linewidth-6.00pt\relax}>{\raggedright\arraybackslash}p{\dimexpr 0.2400\linewidth-6.00pt\relax}>{\raggedright\arraybackslash}p{\dimexpr 0.2300\linewidth-6.00pt\relax}>{\raggedright\arraybackslash}p{\dimexpr 0.3000\linewidth-6.00pt\relax}@{}}
\toprule
\textbf{SAA} & \textbf{VAA} & \textbf{위 식의 OCRT RAA} & \textbf{의미} \\
\midrule
\endfirsthead
\toprule
\textbf{SAA} & \textbf{VAA} & \textbf{위 식의 OCRT RAA} & \textbf{의미} \\
\midrule
\endhead
\bottomrule
\endfoot
30\textdegree{} & 30\textdegree{} & 0\textdegree{} & 태양 look 방위와 센서 look 방위 일치 \\
30\textdegree{} & 210\textdegree{} & 180\textdegree{} & 글린트 가지; 정반사 중심에는 SZA=VZA도 필요 \\
350\textdegree{} & 10\textdegree{} & 20\textdegree{} & 0\textdegree{} 경계를 넘는 차이 \\
10\textdegree{} & 350\textdegree{} & 340\textdegree{} & 위 행의 거울 방위 \\
\end{longtable}
}

이 식은 위에 선언한 좌표에서 유도한 변환이며, 특정 위성 제품의 Stokes 기준면까지 자동으로 맞춘다는 뜻은 아니다. 제품의 양의 방위 회전 방향을 반대로 선택하면 \texttt{RAA}가 \texttt{360-RAA}로 바뀌어 \texttt{I/\allowbreak{}Q}는 같아도 \texttt{U} 부호는 뒤집힌다. 편광 자료에는 원래 SAA/VAA, 변환식, 기준면·회전 방향을 함께 기록하고, 주평면 밖의 한 기하에서 signed \texttt{U}를 대조한다. 제품의 방위·Stokes 정의를 모르면 \texttt{I/\allowbreak{}Q}의 값이 맞는 것만으로 \texttt{U} 변환까지 검증했다고 판단할 수 없다.

그림~\ref{fig:geometry-azimuth}는 SAA에서 VAA까지의 회전을 표시한다. 오른쪽 그림은 같은 좌표를 태양 look 방위 기준으로 돌린 것이며, 거울 방위 두 개를 구분해 저장해야 하는 이유도 보여준다.

% Clockwise is an explicitly declared illustrative coordinate convention.
\begin{figure}[H]
\centering
\begin{tikzpicture}[x=1cm,y=1cm,font=\sffamily\fontsize{9}{11}\selectfont,
  ray/.style={line width=1pt,-{Latex[length=2mm]}},
  arcarr/.style={line width=.85pt,-{Latex[length=1.7mm]}},
  note/.style={align=center,text width=6.9cm}]
\path[use as bounding box] (0,-3.1) rectangle (15.8,3.7);
\begin{scope}[shift={(3.6,0)}]
  \node[anchor=north,font=\sffamily\bfseries] at (0,3.7)
    {\FigLang{(a) 위에서 본 SAA, VAA, RAA}{(a) SAA, VAA, RAA viewed from above}};
  \draw[ocrtLine] (0,0) circle (2.05);
  \draw[ocrtGray,dashed,-{Latex[length=2mm]}] (0,0)--(0,2.65);
  \node[anchor=south,align=center] at (0,2.68)
    {\FigLang{공통 방위 영점}{Shared azimuth zero}};
  \draw[ray,orange!85!black] (0,0)--(60:2.45);
  \draw[ray,ocrtTeal] (0,0)--(-30:2.45);
  \node[anchor=west,align=left] at (1.4,2.15)
    {\FigLang{태양 look}{Solar look}\\SAA $=30^\circ$};
  \node[anchor=north,align=center] at (2.0,-1.47)
    {\FigLang{센서 look}{Sensor look}\\VAA $=120^\circ$};
  \draw[arcarr,orange!85!black] (90:.7) arc[start angle=90,end angle=60,radius=.7];
  \node[anchor=east,orange!85!black] at (-.16,.95) {SAA};
  \draw[arcarr,ocrtTeal] (90:1.03) arc[start angle=90,end angle=-30,radius=1.03];
  \node[ocrtTeal,fill=white,inner sep=1pt] at (.96,.75) {VAA};
  \draw[arcarr,ocrtNavy] (60:1.65) arc[start angle=60,end angle=-30,radius=1.65];
  \node[ocrtNavy,fill=white,inner sep=1.5pt] at (2.14,.24) {RAA $=90^\circ$};
  \fill[ocrtNavy] (0,0) circle (2pt);
  \node[anchor=north east] at (-.1,-.1) {\FigLang{화소}{Pixel}};
  \node[note,anchor=north] at (0,-2.28)
    {$\mathrm{RAA}=(\mathrm{VAA}-\mathrm{SAA})\bmod 360^\circ$};
  \node[note,anchor=north] at (0,-2.83)
    {\FigLang{이 그림의 양의 회전: 시계 방향}{Positive rotation in this figure: clockwise}};
\end{scope}
\draw[ocrtLine] (7.8,-2.95)--(7.8,3.35);
\begin{scope}[shift={(11.8,0)}]
  \node[anchor=north,font=\sffamily\bfseries] at (0,3.7)
    {\FigLang{(b) 태양 look을 위쪽으로 맞춘 좌표}{(b) Rotate the solar look to the top}};
  \draw[ocrtLine] (0,0) circle (1.72);
  \draw[ray,orange!85!black] (0,0)--(0,2.02);
  \draw[ray,ocrtTeal] (0,0)--(1.92,0);
  \draw[ray,ocrtTeal] (0,0)--(0,-1.87);
  \draw[ray,ocrtTeal] (0,0)--(-1.92,0);
  \draw[purple!70!black,dashed,line width=.8pt] (0,0)--(70:1.72);
  \draw[purple!70!black,dashed,line width=.8pt] (0,0)--(110:1.72);
  \node[purple!70!black,anchor=west] at (.7,1.69) {$20^\circ$};
  \node[purple!70!black,anchor=east] at (-.7,1.69) {$340^\circ$};
  \node[anchor=south,align=center] at (0,2.1)
    {$0^\circ\;(=360^\circ)$\\\FigLang{태양 look}{Solar look}};
  \node[anchor=west] at (2.05,0) {$90^\circ$};
  \node[anchor=east] at (-2.05,0) {$270^\circ$};
  \node[anchor=north,align=center] at (0,-1.98)
    {$180^\circ$\quad\FigLang{직접 글린트 가지}{Direct-glint branch}};
  \fill[ocrtNavy] (0,0) circle (2pt);
  \node[note,anchor=north] at (0,-2.6)
    {\FigLang{거울 방위 $20^\circ\leftrightarrow340^\circ$: 같은 $I,Q$, 반대 $U$}{Mirror pair $20^\circ\leftrightarrow340^\circ$: same $I,Q$, opposite $U$}};
\end{scope}
\end{tikzpicture}
\caption{\FigLang{위에서 본 수평 look 방위의 차이로 RAA를 정한다. 영점과 시계 방향은 이 그림에서 선택한 규약이다. 제품의 SAA/VAA를 같은 영점·회전 방향으로 맞춘 후 변환한다. 축대칭 조건에서 $20^\circ$와 $340^\circ$를 모두 $20^\circ$로 접으면 signed $U$에 필요한 거울 방향 정보가 사라진다. 실제 편광 비교에는 Stokes 기준면도 맞춘다.}{RAA is the difference between horizontal look azimuths viewed from above. The zero reference and clockwise rotation are choices declared for this figure. Align the product's SAA/VAA origins and rotation directions before conversion. Under axisymmetric conditions, folding both $20^\circ$ and $340^\circ$ to $20^\circ$ discards the mirror-direction information needed for signed $U$. Polarization comparisons also require aligned Stokes reference planes.}}
\label{fig:geometry-azimuth}
\end{figure}


\subsection{0--180\textdegree{}로 접힌 RAA만 있는 경우}
\label{sec:auto:01_concepts_geometry:4.2}
접힌 값은 흔히 다음처럼 만든다.


\begin{ManualCode}
raa_360 = (vaa_deg - saa_deg) % 360.0
raa_folded = min(raa_360, 360.0 - raa_360)
\end{ManualCode}

축대칭 조건에서 \texttt{I}와 \texttt{Q}만 사용하는 LUT에는 반원 격자를 이용할 수 있다. 그러나 \texttt{raa\_\allowbreak{}folded=\allowbreak{}20\textdegree{}}만 남기면 원래 \texttt{20\textdegree{}}였는지 \texttt{340\textdegree{}}였는지 알 수 없다. \texttt{U}를 보존하려면 원래 0--360\textdegree{} 방향 또는 거울 방향 여부를 함께 저장해야 한다. 절댓값 \texttt{abs(U)}를 저장하면 방향 정보가 사라진다.

\section{다른 코드의 Phi와 맞추는 방법}
\label{sec:auto:01_concepts_geometry:5}
현재 \texttt{rt\_\allowbreak{}raa\_\allowbreak{}convention.\allowbreak{}h}는 OCRT--OSOAA 비교에 다음 두 \textbf{동등한 대안 중 하나}를 정의한다.


{\TableFont
\begin{longtable}{@{}>{\raggedright\arraybackslash}p{\dimexpr 0.3000\linewidth-5.33pt\relax}>{\raggedright\arraybackslash}p{\dimexpr 0.2300\linewidth-5.33pt\relax}>{\raggedright\arraybackslash}p{\dimexpr 0.4700\linewidth-5.33pt\relax}@{}}
\toprule
\textbf{변환 방식} & \textbf{OSOAA에 전달할 방위} & \textbf{OCRT와 비교할 U} \\
\midrule
\endfirsthead
\toprule
\textbf{변환 방식} & \textbf{OSOAA에 전달할 방위} & \textbf{OCRT와 비교할 U} \\
\midrule
\endhead
\bottomrule
\endfoot
A & \texttt{Phi\_\allowbreak{}OSOAA=\allowbreak{}(180-RAA\_\allowbreak{}OCRT) mod 360} & \texttt{U\_\allowbreak{}OCRT=\allowbreak{}-U\_\allowbreak{}OSOAA} \\
B & \texttt{Phi\_\allowbreak{}OSOAA=\allowbreak{}(RAA\_\allowbreak{}OCRT+180) mod 36\allowbreak{}0} & \texttt{U\_\allowbreak{}OCRT=\allowbreak{}U\_\allowbreak{}OSOAA} \\
\end{longtable}
}

\texttt{I/\allowbreak{}Q}는 두 방식에서 같은 물리 비교를 구성한다. A 방식의 U 반전과 B 방식의 각도 변환을 동시에 적용하지 않는다. 이 규칙은 해당 비교 하네스의 출력 기준을 전제로 하며, 임의 코드 전체에 적용하는 보편 공식이 아니다.

현재 OCRT의 대기 출력과 수면 위 \texttt{Rrs}, 수면 아래 \texttt{rrs}는 \textbf{같은 공개 RAA}로 Fourier 재구성한다. 과거 보고서의 \texttt{Phi\_\allowbreak{}OSOAA,\allowbreak{}water=\allowbreak{}RAA\_\allowbreak{}OCRT} 무변환 규칙은 수정 전 물 출력에 대응했던 규칙이다. 현재 소스의 물 출력에는 사용하지 않는다.

\texttt{-\allowbreak{}-\allowbreak{}ccrr-\allowbreak{}compare}는 CCRR 참조 CSV를 읽을 때 내부에서 \texttt{RAA\_\allowbreak{}OCRT=\allowbreak{}(180-RAA\_\allowbreak{}CCRR) mod 360}을 적용한다. CSV의 \texttt{raa\_\allowbreak{}deg}가 이미 OCRT 규약이면 이 비교 모드에 그대로 넣어 이중 변환하지 않는다. 결과의 \texttt{raa\_\allowbreak{}ccrr\_\allowbreak{}deg}, \texttt{raa\_\allowbreak{}ocrt\_\allowbreak{}deg}, \texttt{raa\_\allowbreak{}mapping}을 확인한다.

\section{Stokes 부호와 주평면 밖 검증}
\label{sec:auto:01_concepts_geometry:6}
공개 기본 경로는 관측 방향의 자오면(관측 방향과 수직선이 이루는 면)에 맞춘 Stokes 성분을 사용한다. 표면 기본 Q 규약은 \texttt{q\_\allowbreak{}convention=\allowbreak{}1}, Fresnel 핵심 성분 \texttt{Q\_\allowbreak{}kernel=\allowbreak{}(Rp-Rs)/\allowbreak{}2}다. \texttt{Q}와 \texttt{U}는 부호가 있는 값이며 음수도 정상이다. 센서 좌표나 산란면 좌표와 비교할 때는 기준면 회전이 필요할 수 있다.

축대칭 표면·매질 및 같은 기준면에서 다음 대칭성이 성립한다.


\begin{ManualCode}
I(RAA) =  I(360−RAA)
Q(RAA) =  Q(360−RAA)
U(RAA) = −U(360−RAA)
\end{ManualCode}

주평면 \texttt{RAA=\allowbreak{}0/\allowbreak{}180\textdegree{}}에서는 U가 0 또는 수치 오차 수준이라 U 부호 오류가 드러나지 않는다. \texttt{RAA=\allowbreak{}45/\allowbreak{}315\textdegree{}}처럼 주평면 밖의 쌍도 확인한다. 정확히 \texttt{VZA=\allowbreak{}0\textdegree{}}에서는 자오면이 기하적으로 퇴화하므로 RAA별 Q/U가 직관적인 좌우 방향과 대응하지 않을 수 있다. 모델 비교에는 천저값만 쓰지 말고 작은 양의 VZA도 포함한다.

\section{산란각과 과거 주석의 충돌}
\label{sec:auto:01_concepts_geometry:7}
위에서 선언한 look 방향으로부터 실제 태양광 진행 벡터와 상향 관측 벡터를 내적하면 대기 산란각은 다음과 같다.


\begin{ManualCode}
mu_s = cos(SZA), mu_v = cos(VZA)
cos(Theta) = −mu_s*mu_v − sin(SZA)*sin(VZA)*cos(RAA_OCRT)
\end{ManualCode}

이 식은 현재 \texttt{rt\_\allowbreak{}ipss\_\allowbreak{}cos\_\allowbreak{}scattering\_\allowbreak{}angle()}의 식과 일치한다. \texttt{SZA=\allowbreak{}VZA=\allowbreak{}30\textdegree{}}이면 \texttt{RAA=\allowbreak{}0\textdegree{}}에서 \texttt{Theta=\allowbreak{}180\textdegree{}}이고, 글린트 가지인 \texttt{RAA=\allowbreak{}180\textdegree{}}에서는 \texttt{Theta=\allowbreak{}120\textdegree{}}다. \textbf{수면 정반사 방향과 대기 입자의 정확 후방산란 방향은 같은 개념이 아니다.}

현행 저장소에는 해결되지 않은 내부 설명 충돌도 있다. \texttt{rt\_\allowbreak{}raa\_\allowbreak{}convention.\allowbreak{}h}의 \texttt{rt\_\allowbreak{}atm\_\allowbreak{}scatter\_\allowbreak{}cos\_\allowbreak{}from\_\allowbreak{}public\_\allowbreak{}raa()}는 위 식과 반대인 \texttt{+sin(SZA)*sin(VZ\allowbreak{}A)*cos(RAA)}를 사용하고, 관련 주석과 \texttt{test\_\allowbreak{}raa\_\allowbreak{}convention.\allowbreak{}c}는 \texttt{RAA=\allowbreak{}180\textdegree{}}를 정확 후방산란이라고 표현한다. 이번 문서 작성 시점에 이 helper의 생산 \texttt{src} 호출은 확인되지 않았으며, 실제 Fourier 재구성은 \texttt{public\_\allowbreak{}RAA+\ensuremath{\pi}}, IPSS는 위의 \texttt{-cos(RAA)} 식을 사용한다. 따라서 그 helper·과거 감사 보고서의 “정확 후방산란” 문구를 위성 기하 변환의 근거로 사용하지 않는다. 이 매뉴얼은 소스 코드를 변경하지 않는다.

\subsection{구현 확인 위치}
\label{sec:auto:01_concepts_geometry:7.1}
\begin{itemize}
\item \href{https://github.com/brtnt/OCRT/blob/54861c68e35ec5aaa3938fad53e93dc2dfa7eb1d/code/OCRT_C/src/rt_raa_convention.h}{공개 RAA helper 및 최신 OSOAA 매핑}
\item \href{https://github.com/brtnt/OCRT/blob/54861c68e35ec5aaa3938fad53e93dc2dfa7eb1d/code/OCRT_C/src/rt_fourier.c}{Fourier 재구성의 내부 +\ensuremath{\pi} 변환}
\item \href{https://github.com/brtnt/OCRT/blob/54861c68e35ec5aaa3938fad53e93dc2dfa7eb1d/code/OCRT_C/src/rt_ipss.c}{IPSS 산란각·화소 기준 기하}
\item \href{https://github.com/brtnt/OCRT/blob/54861c68e35ec5aaa3938fad53e93dc2dfa7eb1d/code/OCRT_C/src/main.c}{CLI 모드 선택과 CCRR 변환}
\item \href{https://github.com/brtnt/OCRT/blob/54861c68e35ec5aaa3938fad53e93dc2dfa7eb1d/code/OCRT_C/src/shared/surface.c}{표면 Stokes 및 자오면 회전}
\end{itemize}

입력 파라미터를 바꾸기 전에 \texttt{RAA=\allowbreak{}180\textdegree{}} 글린트, \texttt{0\textdegree{}} 반대 가지, 주평면 밖 U 부호라는 세 가지를 먼저 맞추면 비교에서 생기는 대표적인 180\textdegree{}·부호 오류를 줄일 수 있다.

